\centerline{\bf \Huge Алгебраическая диагностика}

\begin{flushright}
        \scriptsize{Cucu-Cucurella se come una paella}
\end{flushright}

\begin{enumerate}

\item Пусть $x,y,z$ — действительные числа, отличные от $1$, причём
\[
    xyz=1.
\]
Докажите, что
\[
    \frac{x^2}{(x-1)^2}
    +\frac{y^2}{(y-1)^2}
    +\frac{z^2}{(z-1)^2}
    \geqslant 1.
\]
% Источник: IMO 2008, задача 2.
% https://mathus.ru/math/doner.pdf

\item Положительные числа $x,y,z$ удовлетворяют неравенству
\[
    xyz\geqslant xy+yz+zx.
\]
Докажите, что
\[
    \sqrt{xyz}\geqslant \sqrt{x}+\sqrt{y}+\sqrt{z}.
\]
% Источник: региональный этап ВсОШ, 2016, 11 класс, задача 2.
% https://olympiads.mccme.ru/vmo/2016/iii-day1.pdf

\item Зафиксировано положительное число $k$. Последовательность
$(x_n)$ задана условиями
\[
    x_1=1,
    \qquad
    x_{n+1}=\frac12\left(x_n+\frac{k}{x_n}\right).
\]
Докажите, что последовательность $(x_n)$ сходится, и найдите её
предел.
% Источник: Н.Б. Алфутова, А.В. Устинов,
% «Алгебра и теория чисел», задача 09.048.
% https://problems.ru/view_problem_details_new.php?id=61299

\item Многочлен $P$ степени $n>1$ имеет $n$ различных действительных
корней $x_1,x_2,\ldots,x_n$. Корни его производной равны
$y_1,y_2,\ldots,y_{n-1}$. Докажите, что
\[
    \frac{x_1^2+x_2^2+\cdots+x_n^2}{n}
    >
    \frac{y_1^2+y_2^2+\cdots+y_{n-1}^2}{n-1}.
\]
% Источник: Турнир городов, 2007/08,
% весенний сложный вариант, старшая группа, задача 3.
% https://problems.ru/view_problem_details_new.php?id=64613

\item Функция $f\colon\mathbb R\to\mathbb R$ удовлетворяет равенству
\[
    (x-1)f\left(\frac{x+1}{x-1}\right)-f(x)=x
\]
при всех $x\ne 1$. Найдите все такие функции.
% Источник: ВсОШ, окружной этап, 1994, 11 класс, задача 6.
% https://mathus.ru/math/funcura.pdf

\item Все коэффициенты ненулевого многочлена принадлежат множеству
$\{-1,0,1\}$. Докажите, что каждый его действительный корень
принадлежит отрезку
\[
    [-2,2].
\]
% Источник: Московская математическая олимпиада, 1965,
% 11 класс, задача 1.
% https://mathus.ru/math/polynom.pdf

\item Докажите, что уравнение
\[
    x^3+y^3=4\bigl(x^2y+xy^2+1\bigr)
\]
не имеет решений в целых числах.
% Источник: ВсОШ, окружной этап, 1993,
% 9--10 классы, задача 5.
% https://mathus.ru/math/ez.pdf

\item Найдите все простые числа $p$, для которых существуют
натуральные числа $x$ и $y$, удовлетворяющие равенству
\[
    p^x=y^3+1.
\]
% Источник: ВсОШ, окружной этап, 2003, 11 класс, задача 1.
% https://mathus.ru/math/ez.pdf

\end{enumerate}